\documentclass[12pt,a4paper]{article}

\usepackage[utf8]{inputenc}
\usepackage[T1]{fontenc}
\usepackage[spanish]{babel}
\usepackage{lmodern}
\usepackage{microtype}
\usepackage{geometry}
\usepackage{setspace}
\usepackage{parskip}

\geometry{
  top=2.5cm,
  bottom=2.5cm,
  left=3cm,
  right=3cm
}

\setstretch{1.15}

\title{De la profunda desigualdad}
\author{Vicente Domínguez Arca}
\date{}

\begin{document}

\maketitle

A menudo, cuanto más esfuerzo pongo en observar las generalidades de las sociedades que llamamos ``occidentales'', termino encontrando, atónito, una filigrana perfectamente común.

Mientras un ensordecedor clamor reclama igualdad de condiciones y posibilidades, igualdad en la contribución al erario, igualdad ante las instituciones e iguales servicios, ese mismo clamor, como si de una suerte de ambivalencia cuántica se tratara, reclama simultáneamente la libertad del individuo para acumular.

No se trata necesariamente de una paradoja.

Tampoco de una simple incongruencia.

Debe tratarse de algo mucho más fundamental, más original, más radical.

Digo yo que debe tratarse de un problema.

Y, a mi entender, no tiene una solución trivial.

Más bien se antoja complicada.

Propongo un ejercicio mental.

Consiste en abstraernos por un momento de la realidad concreta para intentar determinar cuál de esas dos perspectivas podría ser prístina, si acaso lo es alguna.

Sin necesidad de traer a Rousseau al presente, pensemos en la posibilidad de que el ser humano sea originalmente un ser cuyo horizonte procura el bienestar social, entendido no como una igualdad matemática entre todos los individuos, sino como la conservación de unas condiciones que permitan la existencia y continuidad del tejido del que forma parte.

Supongamos que esto es así.

Si el ser humano es radicalmente social, su existencia individual no puede comprenderse completamente al margen de los demás. Nace de otros, aprende de otros, recibe un lenguaje, conocimientos, instituciones, cuidados y posibilidades construidas antes de su propia existencia.

El individuo aislado puede ser pensado como abstracción.

Resulta mucho más difícil encontrarlo en su realidad natural.

Estamos marcados por la necesidad del prójimo.

Somos nodos de un tejido que nos precede, que contribuimos a transformar y que probablemente continuará después de nosotros.

Si este carácter social es originario, parece razonable pensar que existe también un horizonte de conservación de las condiciones que permiten la continuidad de ese tejido.

No hablo necesariamente de altruismo.

Tampoco de bondad.

Mucho menos de una supuesta inclinación natural hacia una igualdad absoluta.

Hablo de algo más elemental: si nuestra existencia depende radicalmente de una estructura social, la conservación de unas condiciones mínimas de integridad de esa estructura forma parte de las condiciones de nuestra propia existencia.

Ahora introduzcamos la acumulación.

Conviene precisar qué entiendo aquí por ella.

No me refiero al simple acopio de recursos.

Guardar alimento para mañana, conservar herramientas, construir un refugio o disponer de determinados bienes capaces de proporcionar seguridad futura no constituye necesariamente el fenómeno que intento señalar.

La acumulación que aquí interesa comienza cuando la apropiación diferencial supera aquella función inmediata y genera una asimetría creciente dentro del tejido; especialmente cuando aquello acumulado permite, además, aumentar la capacidad futura de seguir acumulando.

En ese momento aparece una dinámica particular.

La acumulación produce capacidad de acumulación.

La desigualdad produce capacidad de aumentar la desigualdad.

Y aquello que inicialmente podría contemplarse como una diferencia cuantitativa entre nodos comienza a convertirse en un gradiente dentro del propio tejido.

Si esto es así, carecería de sentido afirmar sin más que la tendencia hacia esa acumulación y la tendencia hacia la conservación del bienestar social son dos expresiones igualmente originarias de una misma sociabilidad.

Aparece una suerte de problema de los contrarios.

Por un lado, una tendencia hacia la conservación del tejido.

Por otro, una tendencia cuya realización ilimitada genera asimetrías capaces de tensionarlo y, finalmente, deteriorarlo.

No se trata aquí de diferentes personalidades.

Naturalmente existirán personas más generosas, más egoístas, más prudentes, más ambiciosas o más indiferentes.

Eso no es lo que intento plantear.

Hablo de rasgos sociales considerados como principios originarios.

Si una determinada tendencia tiene como horizonte la conservación de las condiciones que permiten la existencia del tejido, resulta difícil considerar igualmente originaria otra tendencia cuyo desarrollo conduce precisamente a erosionar esas condiciones.

Hagamos ahora el ejercicio contrario.

Y, sin necesidad de traer a nuestros días a Hobbes, supongamos que el ser humano procura prístinamente la acumulación individual, buscando con ella asegurar su propia supervivencia, la de su estirpe o la ampliación permanente de su capacidad sobre los recursos disponibles.

La imagen resultante es diferente.

Si cada nodo procura aumentar indefinidamente su posición relativa mediante la acumulación, el tejido social queda sometido a una tensión permanente.

Cada incremento diferencial modifica la relación con los demás.

Y si los recursos, las posibilidades o las capacidades de apropiación no son ilimitados, la acumulación de unos terminará afectando, directa o indirectamente, a la disponibilidad de otros.

La sociedad aparecería entonces como un entramado de individuos permanentemente enfrentados, no necesariamente mediante violencia explícita, pero sí mediante una lucha constante por posiciones, recursos y capacidad futura de apropiación.

Ahora tenemos las dos imágenes.

¿Qué sentido tienen ambas?

¿Pueden constituir simultáneamente dos horizontes originarios de nuestra condición social?

Este autor encuentra grandes dificultades para aceptarlo.

No porque un organismo no pueda albergar impulsos contradictorios.

Puede hacerlo.

No porque una persona no pueda experimentar simultáneamente deseos de cooperación y competencia.

Evidentemente puede.

El problema aparece cuando pretendemos elevar ambas dinámicas a principios originarios de aquello que nos constituye socialmente.

Si nuestra condición humana es radicalmente relacional y la existencia del individuo depende del tejido, una dinámica cuya realización ilimitada deteriora sistemáticamente las condiciones de ese mismo tejido difícilmente puede constituir, al mismo nivel, el principio que explica nuestra sociabilidad.

Por eso este autor considera plausible otra posibilidad.

Que el deseo de acumulación, al menos en la forma social que aquí estamos considerando, no sea originario, sino impuesto.

Pero debo detenerme inmediatamente en esta palabra.

Por \emph{impuesto} no entiendo necesariamente coaccionado.

No estoy imaginando una fuerza exterior que obligue violentamente al individuo a acumular.

Utilizo el término en un sentido más radical: aquello que, no perteneciendo originariamente al individuo, llega a incorporarse a él mediante su interacción con el entorno.

Puede ser aprendido.

Puede ser imitado.

Puede ser incentivado.

Puede surgir de determinadas necesidades.

Puede construirse culturalmente.

Puede aparecer como consecuencia de instituciones, costumbres, expectativas o recompensas sociales.

Incluso puede llegar a ser experimentado por el propio individuo como un deseo completamente suyo.

Nada de esto exige violencia.

Basta con que exista tejido.

Porque si somos nodos interconectados, aquello que sucede en el tejido participa también en la configuración de los nodos.

Pero existe una tensión estructural difícil de ignorar.

¿Hasta qué punto puede aumentar la desigualdad derivada de la libertad de acumulación antes de deteriorar las condiciones que permiten continuar considerando iguales a quienes forman parte del mismo tejido?

Y llegados aquí aparece inevitablemente otra pregunta.

Sin querer generar una lectura demasiado simplista del capitalismo, ¿qué papel desempeña el neoliberalismo en esta dinámica?

Quizá su función no consista en haber inventado la acumulación.

Ni siquiera sería necesario atribuirle la creación de un supuesto egoísmo humano.

La cuestión podría ser mucho más sutil.

¿Qué ocurre si una determinada organización económica, política y cultural convierte la acumulación en una expresión privilegiada de la libertad individual?

¿Qué sucede si la capacidad de acumular comienza a interpretarse como resultado fundamentalmente individual, mientras se difuminan las innumerables relaciones sociales que hicieron posible aquella acumulación?

Nadie acumula completamente solo.

Para que exista patrimonio deben existir instituciones que reconozcan su titularidad.

Para que exista intercambio deben existir otros con quienes intercambiar.

Para que exista producción deben existir conocimientos, infraestructuras, recursos y relaciones.

Para que exista riqueza socialmente reconocible debe existir sociedad.

El nodo acumulador nunca deja realmente de depender del tejido.

Puede, sin embargo, llegar a imaginarse desacoplado de él.

Y quizá sea ahí donde debamos buscar una de las claves.

\end{document}